\section{Model Comparison for Sentiment Analysis}

\subsection{Introduction}

This section presents the comparison of various machine learning models trained for sentiment analysis. The models include traditional classifiers such as \textbf{Logistic Regression}, \textbf{Decision Tree}, \textbf{XGBoost}, \textbf{Random Forest}, \textbf{Perceptron}, and \textbf{Naïve Bayes}, along with more complex models like \textbf{CNN-LSTM}, \textbf{HMM}, and \textbf{Bayesian Networks}. Each model's performance is evaluated using key metrics: \textbf{Accuracy}, \textbf{Precision}, \textbf{Recall}, \textbf{F1-score}, and \textbf{ROC AUC}.

\subsubsection{Model Training and Evaluation Workflow}

To ensure a robust evaluation of sentiment classification models, the following steps are undertaken:

\begin{itemize}
    \item \textbf{Instantiating a GridSearch Object:} The model is initialized with a set of hyperparameters using \textit{GridSearchCV}, allowing an exhaustive search over different hyperparameter combinations to identify the optimal settings.
    
    \item \textbf{Fitting the Training Data:} The training dataset is fed into the model, enabling it to learn patterns that distinguish between different sentiment classes.
    
    \item \textbf{K-Fold Cross-Validation:} To enhance generalization, \textit{K-Fold Cross-Validation} is applied, dividing the dataset into multiple subsets to train and validate the model iteratively.
    
    \item \textbf{Saving the Trained Model:} The best-performing model, based on cross-validation results, is stored for future use, ensuring consistency in later inference stages.
    
    \item \textbf{Testing on a Separate Dataset:} The trained model is evaluated on a test dataset to measure its real-world generalization ability, providing a reliable estimate of performance.
    
    \item \textbf{Logging Performance Metrics:} Key metrics such as \textbf{Accuracy}, \textbf{Precision}, \textbf{Recall}, \textbf{F1-score}, and \textbf{ROC AUC} are recorded to facilitate model comparison.
\end{itemize}

This workflow ensures a structured and reliable methodology for training, validating, and comparing sentiment analysis models, aiding in the selection of the most effective approach.

\subsubsection{Evaluation Metrics Overview}

To assess the performance of sentiment analysis models, we utilize five key evaluation metrics: \textbf{Accuracy}, \textbf{Precision}, \textbf{Recall}, \textbf{F1-score}, and \textbf{ROC AUC}. 

\begin{itemize}
    \item \textbf{Accuracy} measures the overall correctness of the model by calculating the proportion of correctly classified instances. However, it may not always be the best metric when dealing with imbalanced sentiment classes. 
    
    \item \textbf{Precision} quantifies the proportion of correctly predicted positive samples out of all predicted positive samples, which is crucial when minimizing false positives, such as in cases where detecting negative sentiment is critical. 
    
    \item Conversely, \textbf{Recall} indicates how well the model identifies actual positive cases, making it essential for scenarios where missing positive sentiment (e.g., detecting customer dissatisfaction) is more detrimental.  
    
    \item The \textbf{F1-score} provides a balanced measure of both precision and recall, ensuring that the model maintains strong predictive power across both metrics. 
    
    \item Lastly, \textbf{ROC AUC} (Receiver Operating Characteristic - Area Under the Curve) evaluates the model's ability to distinguish between different sentiment classes, providing insight into its overall discriminative power. By considering these metrics, we can determine the best-performing model based on different application needs, balancing between false positives and false negatives.
\end{itemize}

\subsection{Baseline Model Comparison (Best-performing Feature Methods)}

The performance of different models for sentiment analysis using the best feature extraction method per model is summarized in Table~\ref{tab:model_comparison}.

\begin{table}[H]
\centering
\caption{Model Comparison for Sentiment Analysis (Best-performing features applied)}
\label{tab:model_comparison}
\renewcommand{\arraystretch}{1.5}
\begin{tabular}{||c|c|c|c|c|c||}
\hline
\textbf{Model} & \textbf{Accuracy} & \textbf{ROC AUC} & \textbf{F1} & \textbf{Precision} & \textbf{Recall} \\
\hline
HMM & 0.5141 & 0.4997 & 0.6697 & 0.5156 & 0.9552 \\
Bayesian Network & 0.6495 & 0.7143 & 0.6875 & 0.6364 & 0.7476 \\
LDA (Word2Vec) & 0.7160 & 0.7846 & 0.7290 & 0.7127 & 0.7461 \\
Logistic Regression (TF-IDF) & \textbf{0.7527} & \textbf{0.8299} & \textbf{0.7681} & 0.7413 & 0.7968 \\
XGB (Word2Vec) & 0.7168 & 0.7146 & 0.7495 & 0.7316 & 0.7942 \\
Random Forest (CountVectorizer) & 0.7300 & 0.8000 & 0.7600 & 0.7000 & 0.8200 \\
BiLSTM with CRF & 0.7092 & 0.7082 & 0.7326 & 0.6853 & 0.7868 \\
Stacking (Selected Models) & 0.7451 & 0.7431 & 0.7657 & 0.7270 & 0.8087 \\
Voting (Multiple Logistic Regressions) & 0.7373 & 0.7356 & 0.7268 & 0.7232 & 0.7938 \\
\hline
\end{tabular}
\end{table}

\textbf{Discussion:} Logistic Regression with TF-IDF emerges as the best performer across metrics, closely followed by ensemble-based models like Random Forest and Stacking. HMM, while achieving high recall, shows poor ROC AUC and overall balance.

\subsection{Discussion on Feature Selection and Dimensionality Reduction Methods}

\subsubsection{Chi-Squared Feature Selection}

As shown in Table~\ref{tab:full_results}, Chi-Squared feature selection enhances the performance of models by selecting features that are most statistically correlated with the sentiment class labels. This univariate method is particularly beneficial for high-dimensional sparse data like TF-IDF or CountVectorizer outputs.

${}$\\
\textbf{Model Analysis:} 

Tree-based models such as \textbf{Random Forest (Accuracy = 0.5986, F1 = 0.6718)} and \textbf{XGBoost (Accuracy = 0.6005, F1 = 0.6686)} exhibit excellent synergy with Chi-Squared features, demonstrating their robustness in handling structured, selected input. \textbf{Perceptron-TFIDF}, on the other hand, shows poor F1 score despite high precision (Precision = 0.6233, F1 = 0.3114), likely due to overfitting on low-frequency selected features.

${}$\\
\textbf{Comparison with Theory:} 

Chi-Squared aligns with theoretical expectations—effective in eliminating irrelevant features and enhancing performance of linear and ensemble models that assume feature independence or separability.

\textbf{Limitations:}

\begin{itemize}
    \item May discard semantically important but statistically rare features.
    \item Ignores feature interaction effects, which can reduce contextual understanding in complex models.
\end{itemize}

\textbf{Use-case Fit:}

\begin{itemize}
    \item Suitable for rapid prototyping or production pipelines using traditional models like Logistic Regression, SVM, or XGBoost.
    \item Less effective in deep or contextual models requiring nuanced semantic relationships.
\end{itemize}

\subsubsection{Topic Modeling Feature Selection}

Table~\ref{tab:topic_modeling_results} shows the application of LDA-based topic modeling for feature selection. This approach attempts to reduce dimensionality while preserving semantic structure by selecting features tied to dominant latent topics.

${}$\\
\textbf{Model Analysis:} 

\textbf{Random Forest (Accuracy = 0.5760, ROC AUC = 0.6091)} and \textbf{MLP (F1 = 0.5912)} show improved contextual robustness with topic-based features, suggesting they can capture topic distribution patterns well. However, linear models such as \textbf{Logistic Regression} and \textbf{SVM} offer moderate gains, while \textbf{Perceptron-TFIDF} lags behind due to underutilization of latent structures.

${}$\\
\textbf{Comparison with Theory:} 

As expected, models capable of learning from dense, semantically-clustered data (e.g., MLP and Random Forest) align well with topic model output. Feature coherence across documents helps reinforce pattern learning.

\textbf{Limitations:}

\begin{itemize}
    \item Quality of topic extraction heavily depends on preprocessing and corpus size.
    \item Some features may be generic across classes, reducing discriminative power.
\end{itemize}

\textbf{Use-case Fit:}

\begin{itemize}
    \item Effective for thematic or long-document sentiment classification.
    \item Less suitable for microtext or domain-specific applications with fine-grained labels.
\end{itemize}

\subsubsection{Variance Threshold Feature Selection}

Also reflected in Table~\ref{tab:full_results}, variance thresholding removes features with minimal value change across documents, assuming that these are less informative.

${}$\\
\textbf{Model Analysis:} 

\textbf{MLP} and \textbf{XGBoost} continue to benefit by focusing only on highly active features, as their architectures adapt well to compressed input spaces. \textbf{Naïve Bayes} shows less consistent performance since it relies on full feature distributions for probabilistic modeling.

${}$\\
\textbf{Comparison with Theory:} 

This method supports the theoretical assumption that low-variance features do not contribute significantly to model decision boundaries. Models that thrive on dense signals (e.g., MLP, XGBoost) are favored.

\textbf{Limitations:}

\begin{itemize}
    \item Eliminates features that may appear rarely but carry strong sentiment signals (e.g., "hate", "never").
    \item Can degrade generative or probabilistic models that depend on count statistics.
\end{itemize}

\textbf{Use-case Fit:}

\begin{itemize}
    \item Ideal for real-time sentiment systems with limited memory or storage budgets.
    \item Not recommended in legal or clinical sentiment tasks where low-frequency terms carry weight.
\end{itemize}

\subsubsection{PCA (Principal Component Analysis)}

Table~\ref{tab:pca-testing} displays results using PCA as a dimensionality reduction technique. PCA transforms feature spaces into orthogonal components ranked by variance contribution.

${}$\\
\textbf{Model Analysis:} 

\textbf{XGBoost-TFIDF (Accuracy = 0.5622, F1 = 0.6275, ROC AUC = 0.5781)} performs best among PCA-enabled models, benefiting from its ensemble structure's resilience to feature transformation. However, most models—especially \textbf{Perceptron (Count)} and \textbf{Logistic Regression}—see reduced F1 scores compared to non-PCA setups.

${}$\\
\textbf{Comparison with Theory:} 

PCA’s decorrelation of features helps suppress noise and redundancy. Yet, transformation into latent components hurts interpretability and weakens model alignment with the original semantic structure of TF-IDF/Count features.

\textbf{Limitations:}

\begin{itemize}
    \item Loss of direct feature meaning limits explainability.
    \item Decreased performance in models relying on lexical-level granularity.
\end{itemize}

\textbf{Use-case Fit:}

\begin{itemize}
    \item Suitable for storage-constrained applications or embedding pipelines.
    \item Not ideal where interpretability or word-level insights are critical (e.g., education, legal).
\end{itemize}

\subsubsection{Raw TF-IDF and CountVectorizer (Without PCA)}

Table~\ref{tab:testing-count-tfidf} reveals the strongest results across all feature methods, using raw uncompressed TF-IDF and CountVectorizer representations.

${}$\\
\textbf{Model Analysis:} 

\textbf{Logistic Regression (Count)} achieves state-of-the-art Accuracy = \textbf{0.7557} and F1 = \textbf{0.7726}, closely followed by \textbf{MLP} and \textbf{XGBoost}. These results reinforce the effectiveness of sparse high-dimensional features for linear and shallow neural models.

${}$\\
\textbf{Comparison with Theory:} 

TF-IDF and CountVectorizer maintain lexical and frequency signals, directly feeding interpretable input into models. This favors methods like Logistic Regression, which capitalize on such token-level sparsity.

\textbf{Limitations:}

\begin{itemize}
    \item High dimensionality increases training time and risk of overfitting.
    \item Less suitable for deployment on low-resource environments.
\end{itemize}

\textbf{Use-case Fit:}

\begin{itemize}
    \item Recommended for applications requiring transparency, accuracy, and speed.
    \item Avoid in low-latency or mobile environments unless dimensionality is reduced.
\end{itemize}

\subsubsection{GloVe Embedding Features}

In Table~\ref{tab:glove_results} (Table 5.17), models trained on GloVe pre-trained embeddings—dense vectors capturing global co-occurrence—show remarkable generalization.

${}$\\
\textbf{Model Analysis:} 

\textbf{XGBoost (Accuracy = 0.7198, F1 = 0.7321, ROC AUC = 0.7939)} and \textbf{Logistic Regression (F1 = 0.7281)} lead in performance. These models benefit from GloVe’s contextual strength, yielding both syntactic and semantic cues.

${}$\\
\textbf{Comparison with Theory:} 

GloVe aligns with modern embedding theory: mapping similar words close in space and preserving co-occurrence patterns. Its effectiveness in general-purpose sentiment tasks is reaffirmed here.

\textbf{Limitations:}

\begin{itemize}
    \item Cannot adapt to domain-specific vocabulary (e.g., technical jargon).
    \item Needs preprocessing alignment (tokenization, padding).
\end{itemize}

\textbf{Use-case Fit:}

\begin{itemize}
    \item Strong choice for multilingual, general-topic sentiment analysis with pretrained knowledge.
    \item Not recommended where domain-specific embeddings (e.g., BioBERT, FinBERT) would yield better results.
\end{itemize}

\subsection{Comparison with Theory}

\begin{itemize}
    \item \textbf{Logistic Regression:} Strong performance with TF-IDF matches its theoretical efficiency on linearly separable data.
    \item \textbf{Random Forest / XGBoost:} Their robustness to feature noise is confirmed under different feature methods.
    \item \textbf{LDA:} Works well with topic-coherent inputs but is weaker than discriminative models.
    \item \textbf{BiLSTM-CRF:} Expected to model sequence well, but limited data may hinder performance.
    \item \textbf{HMM / Bayesian Networks:} These models underperform due to oversimplified assumptions and data sparsity.
    \item \textbf{Ensemble Methods:} Theoretical stability is reflected in consistent, high-performing results.
\end{itemize}

\subsection{Limitation Discussion}

\textbf{General:}

\begin{itemize}
    \item Dataset size and class imbalance impact complex models.
    \item Feature methods vary in ability to capture contextual meaning.
\end{itemize}

\textbf{Model-specific:}

\begin{itemize}
    \item \textbf{HMM, Bayesian:} Struggle with sparse and high-dimensional data.
    \item \textbf{LDA:} Assumes overly simplistic distributions.
    \item \textbf{Logistic Regression:} Can fail on non-linear data.
    \item \textbf{Random Forest / XGBoost:} Need hyperparameter tuning to avoid noise sensitivity.
    \item \textbf{BiLSTM-CRF:} Requires large datasets to avoid overfitting.
    \item \textbf{Stacking/Voting:} Performance depends on base learners; may add complexity.
\end{itemize}

\subsection{Use-case Fit Analysis}

\textbf{Best suited for:}

\begin{itemize}
    \item \textbf{Logistic Regression (TF-IDF):} Fast, interpretable, production-ready.
    \item \textbf{Random Forest / XGBoost:} Strong generalization for nonlinear tasks.
    \item \textbf{Stacking / Voting:} High-stability applications.
    \item \textbf{BiLSTM-CRF:} Context-sensitive tasks with ample data.
\end{itemize}

\textbf{Less suitable for:}

\begin{itemize}
    \item \textbf{HMM, Bayesian Networks:} Weak on modern, sparse textual features.
    \item \textbf{LDA:} Poor for non-topic-focused classification.
\end{itemize}

\newpage
